% UBT-AI-PROVENANCE-BEGIN
% schema: ubt-ai-provenance/v1
% tier: C_working
% ai_assistance: disclosed
% human_review: risk-based
% editorial_responsibility: Ing. David Jaroš
% policy: ../../../AI_PROVENANCE.md
% notice: Working material; exhaustive human review is not claimed.
% UBT-AI-PROVENANCE-END

\chapter{Lorentz geometry on the Hermitian section \(\mathbb H_{\mathbb C}\)}
\label{ch:lorentz-hc}

\section{Classic basis: \(\mathrm{Herm}(2,\mathbb C)\)}
This section is standard spinor/Lorentz mathematics.  It is included because
UBT uses the same algebra through the isomorphism
\[
 \mathbb B=\mathbb C\otimes\mathbb H\cong M_2(\mathbb C).
\]
Let $\sigma_i$ be the Pauli matrices and associate a real four-vector
$x=(x^0,x^1,x^2,x^3)$ with the Hermitian matrix
\[
 X=x^0 I_2+x^i\sigma_i
 =\begin{pmatrix}
 x^0+x^3 & x^1-i x^2\\
 x^1+i x^2 & x^0-x^3
 \end{pmatrix}.
\]
A direct determinant calculation gives
\[
 \boxed{\det X=(x^0)^2-(x^1)^2-(x^2)^2-(x^3)^2.}
\]
Thus the determinant on the Hermitian slice is exactly the Minkowski quadratic
form with signature $(+---)$.

\section{Action \(SL(2,\mathbb C)\)}
For $A\in SL(2,\mathbb C)$ define
\[
 X\longmapsto X'=A X A^\dagger .
\]
Hermiticity is preserved because
\[
 (A X A^\dagger)^\dagger=A X A^\dagger.
\]
The determinant is also preserved:
\[
 \det X'=\det A\,\det X\,\det A^\dagger
 =\det X,
\]
since $\det A=1$ and $\det A^\dagger=\overline{\det A}=1$.
Therefore the induced real linear transformation of the four coefficients
preserves the Minkowski norm.  Restricting continuously to the identity gives a
homomorphism
\[
 SL(2,\mathbb C)\longrightarrow SO^+(1,3).
\]
Both $A$ and $-A$ act identically on $X$, so the kernel contains
$\{\pm I_2\}$.  The standard result is the double cover
\[
 \boxed{SL(2,\mathbb C)/\{\pm I_2\}\cong SO^+(1,3).}
\]
No split-quaternion algebra is required for this construction.

\section{Conversion to biquaternion notation}
Use the canonical matrix representation
\[
 \mathbf e_i\longmapsto i\sigma_i.
\]
Then $\sigma_i\leftrightarrow -i\mathbf e_i$, and the Hermitian matrix above
corresponds to the biquaternion
\[
 X_{\mathbb B}=x^0\mathbf1-i x^i\mathbf e_i.
\]
This is simply another coordinate description of the same
$\mathrm{Herm}(2,\mathbb C)$ slice.

UBT's current tetrad convention often uses instead the Lorentz-real basis
\[
 \mathbf u_0=i\mathbf1,
 \qquad \mathbf u_k=\mathbf e_k,
\]
so a tetrad leg is
\[
 E_\mu=i e_\mu{}^0\mathbf1+e_\mu{}^k\mathbf e_k.
\]
With quaternion conjugation $\sharp$,
\[
 E_\mu^\sharp=i e_\mu{}^0\mathbf1-e_\mu{}^k\mathbf e_k.
\]
Using $\mathbf e_j\mathbf e_k+\mathbf e_k\mathbf e_j=-2\delta_{jk}$,
one obtains
\begin{align*}
 \frac12(E_\mu^\sharp E_\nu+E_\nu^\sharp E_\mu)
 &=\left(-e_\mu{}^0e_\nu{}^0+e_\mu{}^k e_\nu{}^k\right)\mathbf1\\
 &=g_{\mu\nu}\mathbf1.
\end{align*}
This is the same Lorentz geometry with the opposite overall signature
convention $(-+++)$.  A signature convention is bookkeeping; it is not a new
physical degree of freedom.

\section{Why it matters to UBT}
The classical result establishes three useful facts before any UBT dynamics is
introduced:
\begin{enumerate}
 \item complexified quaternions already carry the spinorial Lorentz action;
 \item a four-dimensional Lorentzian quadratic form can be read internally from
       the algebra;
 \item left/right multiplication matters, because the Lorentz action is
       naturally two-sided in matrix language.
\end{enumerate}

What this classical construction does\emph{not} prove is equally important.
It does not prove that the UBT master field dynamically generates every desired
tetrad, it does not derive Einstein's equations, and it does not select the
physical connection.  Those are UBT-specific bridge problems addressed in the
covariant-tetrad chapter and the canonical GR closure ledger.

\section{Specific boost as a control example}
Take a boost of rapidity $\chi$ along the third spatial axis,
\[
 A=\exp\left(\frac{\chi}{2}\sigma_3\right)
 =\begin{pmatrix}e^{\chi/2}&0\\0&e^{-\chi/2}\end{pmatrix}.
\]
Applying $X'=AXA^\dagger$ and comparing coefficients gives
\[
 \begin{pmatrix}x'^0\\x'^3\end{pmatrix}
 =\begin{pmatrix}\cosh\chi&\sinh\chi\\
                  \sinh\chi&\cosh\chi\end{pmatrix}
  \begin{pmatrix}x^0\\x^3\end{pmatrix},
 \qquad x'^1=x^1,\quad x'^2=x^2.
\]
Then
\[
 (x'^0)^2-(x'^3)^2=(x^0)^2-(x^3)^2,
\]
which is the ordinary Lorentz boost written as an internal algebra action.

\section{Living sources and historical note}
The algebra convention is kept in
\texttt{canonical/algebra/biquaternion\_algebra.tex}; current GR use of the
Lorentz-real tetrad is documented in the active GR track.  The former
\path{appendix_P6_lorentz_in_HC.tex} survives in the historical archive,
but the live textbook no longer imports it.
